%% 30_flow.tex -- Chapter 30: Open-Source Design and Verification Flow
%% Status: OUTLINE (build stage 6). Replace the outline with the full chapter when it is written.
\chapter{Open-Source Design and Verification Flow}\label{ch:flow}
\chapterlead{Every circuit in this tutorial can be designed, verified and taped out with free tools and an open PDK.}

\begin{specnote}[Chapter status: outline]
This chapter is an outline. It will be written in build stage 6 (Part V, Lab 14, the capstone and the appendices). The topics, the design question and the labs below are fixed; the derivations, worked examples and figures are still to come.
\end{specnote}

\begin{objectives}
\begin{itemize}
  \item set up IIC-OSIC-TOOLS with Xschem, ngspice/Xyce and KLayout;
  \item write and simulate SAR logic in Verilog with mixed-signal co-simulation;
  \item run corners and Monte Carlo;
  \item run DRC, LVS and parasitic extraction in KLayout;
  \item use the JKU AMS chip template and submit to an IHP OpenMPW shuttle.
\end{itemize}
\end{objectives}

\begin{insight}[Design question]
How do you take an \OB{} block from schematic to a verified GDS that IHP will accept?
\end{insight}

\section{The tool set}\label{sec:ch30-1}
IIC-OSIC-TOOLS~\cite{iicosic}: Xschem, ngspice, Xyce, KLayout, Verilator and Icarus Verilog in one container, with the IHP-Open-PDK~\cite{ihppdk}.

\section{Schematic and simulation}\label{sec:ch30-2}
Xschem testbenches; ngspice and Xyce for AC, noise, transient and transient-noise analyses.

\section{Mixed-signal co-simulation}\label{sec:ch30-3}
Verilog SAR logic with the analog CDAC and comparator.

\section{Corners and Monte Carlo}\label{sec:ch30-4}
Process, voltage and temperature corners; mismatch; what a biomedical chip must survive.

\section{Layout and verification}\label{sec:ch30-5}
KLayout DRC, LVS and PEX; post-layout simulation.

\section{From template to tape-out}\label{sec:ch30-6}
The JKU ihp-sg13g2-ams-chip-template~\cite{jkutemplate} with its padframe, and the IHP OpenMPW route documented in IHP-Open-DesignLib~\cite{ihpdesignlib}.

\begin{tryit}[Labs for this chapter]
Lab~14 (\cref{sec:lab14}) lays out the comparator and CDAC, runs DRC, LVS and PEX with the JKU template and compares post-layout with schematic simulation.
\end{tryit}

\paragraph{Read with this chapter} Dorrer and Pretl~\cite{jkutemplate}, IIC-OSIC-TOOLS~\cite{iicosic}, IHP-Open-DesignLib~\cite{ihpdesignlib}.
